\documentclass[letterpaper,11pt]{moderncv}
\moderncvstyle{banking}
\moderncvcolor{black}

\usepackage[letterpaper, margin=0.8in]{geometry}
\usepackage{xcolor}
\usepackage{fontspec}
\setmainfont{Times New Roman}
\usepackage[unicode]{hyperref}

\name{Saeyoung}{Kim}
\address{Cambridge, MA, USA}{}{}
\phone[mobile]{+1 (351) 322 5510}
\email{saeyoung\_kim@uml.edu}
\social[linkedin]{saeyoung-macx-kim}

\recipient{Hiring Team}{Boston Dynamics\\Waltham, MA}
\date{\today}
\opening{Dear Hiring Team,}
\closing{Sincerely,}

\begin{document}

\makelettertitle

I'm applying for the Staff Systems Test Engineer role (R2380) on the Embodied AI Safety R\&D team. Building the test and validation framework that lets you make real safety claims about robots working alongside people is exactly the kind of problem I want to work on.

Most of my testing background comes from satellite programs at Hanwha Systems, where I led system-level test campaigns for military SAR satellites across EM, QM, and FM builds. That work was structured and evidence-driven by necessity. I developed test plans and acceptance criteria, built EGSE and test stands, executed qualification campaigns, and compiled the requirements traceability and test evidence that went into formal acceptance reviews. When hardware failed during integration, I ran root cause analysis, traced failures through schematics and firmware logs, and coordinated fixes across electrical, mechanical, and software teams. I also provided design-for-test and design-for-safety feedback based on failure trends across build campaigns, and authored the controlled procedures and technical documentation used for both engineering and customer review. The discipline of building a structured evidence base for a complex integrated system is something I'd bring directly to the safety case work on this team.

My PhD was seven years of building and testing integrated sensor systems from scratch, working in a small interdisciplinary team. I designed a multi-modal sensing payload for a CubeSat, covering the full stack from requirements through circuit design, PCB layout, firmware in C/C++, and system-level qualification. I built the test stands, ran environmental and reliability testing across multiple prototype cycles, and used the data to identify failure modes and validate design changes. That experience gave me deep comfort working across hardware and software boundaries to prototype and validate integrated systems.

At UMass Lowell, I've been developing Python tools for statistical process monitoring and automated anomaly detection across 100+ production batches. The core of that work is using data to build evidence-based assessments of system behavior, which maps well to compiling statistical safety evidence for autonomous systems.

I'll be upfront that my safety standards experience is limited. I don't have direct experience with IEC 61508 or ANSI/RIA R15.08, but I've worked within structured qualification frameworks for military satellite hardware that share the same underlying philosophy: define requirements, perform structured analysis, execute tests, and compile evidence to support specific claims about system behavior. I'm confident I can get up to speed on the robotics-specific standards quickly.

I live in Cambridge, so Waltham is local. Regarding work authorization: I am currently authorized to work in the United States and do not require sponsorship. My green card is pending, and I'm happy to discuss details directly.

\makeletterclosing

\end{document}
